August 3, 2026

Editor-in-Chief
Journal of Systems Architecture
Elsevier

SUBJECT: Submission of Manuscript P4 — Spatiotemporal Relational Logic and Kinematic Velocity Filtering for Schedule Compliance

Dear Editor-in-Chief and Associate Editors,

We are pleased to submit our original research manuscript titled "Spatiotemporal Relational Logic and Kinematic Velocity Filtering for Schedule Compliance" for publication consideration in the Journal of Systems Architecture (Elsevier).

1. SUBSYSTEM BACKGROUND AND MOTIVATION
Automated attendance verification systems often suffer from high false alarm rates caused by transient student movements in hallways or camera matching noise. Conventional systems rely on static proximity thresholds without reasoning about spatial context or physical travel constraints. This paper presents a temporal logic framework that correlates localized student observations with institutional schedules.

2. CORE TECHNICAL CONTRIBUTIONS AND NOVELTY
  - Kinematic Teleportation Heuristic: Formulates a spatial velocity filter (v_i <= v_max) that calculates bounding box transition speeds between sequential camera detections, automatically identifying and discarding physically impossible tracking jumps.
  - Spatiotemporal Relational Logic: Integrates institutional timetable functions with temporal debouncing (30-second window) to evaluate attendance compliance and flag truancy anomalies.
  - Empirical Performance: Demonstrates an 85% reduction in false alert generation compared to naive proximity matching, achieving a 98.2% truancy detection F1-score across 52,203 evaluation epochs.

3. SUITABILITY FOR JOURNAL OF SYSTEMS ARCHITECTURE
This manuscript fits the Journal of Systems Architecture's focus on system-level logic design, real-time event reasoning, embedded software architectures, and automated compliance verification.

4. ETHICS AND ORIGINALITY DECLARATIONS
This work is original, has not been published previously, and is not currently under consideration by another journal.

Sincerely,

Polisetti Narendra (Corresponding Author)
Department of Computer Science & Engineering
Swarnandhra College of Engineering and Technology (Autonomous)
Seetharampuram, Narasapur-534280, Andhra Pradesh, India
Email: narendraa1996@gmail.com
